\section{A Two-Earner Model of Undiversified Family Insurance}\label{sec:model}

\providecommand{\E}{\mathbb{E}}

The evidence identifies the pieces: a joint job loss that is a same-month coincidence carried by shared workplaces, a
smaller permanent part shared by couples who hold risky jobs, a wife who searches at once and finds work slowly, a
marriage that is strained, consumption insured against one job loss and perhaps not against two, and a safety net
whose programs pay in different states. A model is needed to put them together, to measure what the correlation costs
the household in the months of the shock, and to evaluate policies the data cannot vary. This section sets out the
model to be estimated; the estimates, the fit and the counterfactuals follow in the next version of the paper.

The period is a month: the joint loss is a coincidence within the month (Figure~\ref{fig:cps_hazard}), the wife's
response builds up over about six months (Section~\ref{sec:household}), unemployment insurance lasts 26 weeks, and
most of the earnings risk of US workers is month-to-month and masked in annual data \citep{ganong2026earnings}. In an
annual period the common and the permanent components of the risk collapse into one same-year correlation. The
closest model is the monthly life-cycle model of joint search of \citet{bacher2025joint}, whose labor-market
block---employment and non-employment with and without unemployment insurance that expires at a constant hazard, a
human capital level for each spouse that accumulates at work and depreciates out of it, savings, and extreme-value
taste shocks solved with the endogenous grid method---is followed here. The model departs from theirs in three ways.
The spouses' displacement risk has a common and a permanent couple component, where their separation shocks are
independent across spouses, which bounds the ratio of the wife's job-loss probability in her husband's loss month to
other months at about 1.25 against about 10 in the CPS.\footnote{With independent hazards that depend only on the
human capital level, the ratio is $1 + \mathrm{cov}(\delta_w, \delta_h)/(\E\delta_w\,\E\delta_h)$; with their hazards
between 1.0\% and 5.0\% a month and a mean near 2\%, perfectly correlated levels give at most about 1.25.} The search
friction is exogenous and race-specific, where their job-finding rates are set by firms. And the government is a set
of schedules---a progressive tax, the Earned Income Tax Credit, food stamps and a consumption floor besides
unemployment insurance---the couple has children whose ages set the cost of the wife's work, the marriage can end, and
couples differ by race. All endogenous parameters are estimated jointly by the simulated method of moments, where
theirs are calibrated.

\subsection{Environment and timing}

\paragraph{Households and time.} The economy is populated by married couples. A couple enters the model in January of
the first calendar year in which the corresponding PSID couple is observed and is followed until the wife is 60: with
the wife aged 25 at the earliest, $t = 0, \dots, 419$ indexes months, and $t$ is a multiple of 12 in every January, so
model years are calendar years. Money is in 2000 dollars a month. The couple is a unitary household that maximises one
utility function in which the spouses share consumption and each spouse's work enters as a disutility; the joint-loss
moments carry no information on a Pareto weight. After a separation the wife is the only decision maker of her
household and the husband leaves the model.

\paragraph{Heterogeneity.} Couples differ in four permanent characteristics, all known to the couple. Race
$r \in \{B, W\}$ enters only through the environment---the displacement hazards, the search friction, the separation
hazard, the initial conditions and the inputs---never through preferences, because the joint loss is the same in both
races while the environment differs (Section~\ref{sec:jointloss}); common preferences make the racial differences in
the wife's response a test of the environment. The workplace match $s \in \{0, 1, 2\}$ sets the common component of
displacement risk: $s = 2$ if the spouses share a detailed industry, a detailed occupation and a class of worker (the
closest the public data come to a shared employer), $s = 1$ if they share a detailed industry but not that cell, and
$s = 0$ otherwise.\footnote{CPS: the wife's same-month excess is 0.337 (s.e.\ 0.015) in the shared-employer cell
(6.3\% of couples, 35\% of the raw excess), 0.211 (s.e.\ 0.009) in the same detailed industry (17.7\%, 53\%) and 0.026
(s.e.\ 0.001) outside the shared-employer cell; if 10--12\% of dual-earner couples work at the same establishment
\citep{hyatt2019coworking}, a shared employer carries a third to a half of the excess. The PSID couple's $s$ at entry
is coded with the CPS definition (\texttt{IGE/60.PaperATargets.py}; CPS: \texttt{IGE/37.CPSJointLossChecks.py}).} A
couple risk type $\tau \in \{L, H\}$ scales both spouses' own hazards and stands for the assortative part of the risk:
wives of men who will be displaced four to six years later are already more often displaced
(Table~\ref{tab:joint_job_loss}), and the CPS elevation in the months around his loss is flat, as a permanent type
predicts. Marital status changes once at most. Education does not index the solution: it enters through each spouse's
initial human capital level and initial wealth.

\paragraph{Choices and timing.} Each month the household chooses consumption and saving subject to a borrowing limit,
the wife's hours $\ell_w \in \{0, P, F\}$ when she has a job (zero means she leaves it), the husband's hours
$\ell_h \in \{P, F\}$ when he has a job, and whether to take up food stamps when eligible. There is no search-effort
choice: a non-employed spouse finds a job with a probability that measures the search friction, and a wife who finds
one decides whether to work.\footnote{``Looking for work'' is therefore a measurement of the model, not a choice: a
non-employed wife is classified as looking when she would accept a job in her current state with probability above one
half; the CPS jump in looking the month after his loss (8.8 points) validates the mapping. The PSID's annual
interview-week status cannot discipline a search choice: a year after the job loss looking is unchanged ($-0.001$,
s.e.\ 0.008; \texttt{IGE/19.FatherDisplacement.py}).} The husband works whenever he has a job.\footnote{Married men
employed the month before are still employed in 99.5\% of cases a month later when they have no involuntary separation
(SIPP; \texttt{IGE/40.SIPPJointLoss.py}); their exits outside displacement are left out.} Within the month: in January
the transitory wage shocks are drawn and observed, and held for the year; taste shocks are drawn for every feasible
alternative $d = (\ell_w, \ell_h, S)$; the household chooses $d$, then consumption and saving; income is received and
spent. At the end of the month the labor-market shocks arrive in this order: the common displacement shock and the own
hazards for spouses who worked, job finding for spouses who were non-employed (an entitled wife then decides whether to
accept the job), and the quit of a wife who chose
$\ell_w = 0$ while holding a job (without entitlement to unemployment insurance). Human capital then moves, the
entitlement states update, the youngest child's age bucket moves, and the separation shock is drawn. A displacement at
the end of month $t$ is a job loss in month $t$ in the data's sense (employed in $t$, not employed in $t+1$), so a
common shock produces both spouses' job losses in the same CPS month; a displaced spouse cannot find a job in the
month of the displacement, as in \citet{bacher2025joint}.

\subsection{Preferences}

The married household's flow utility is
\begin{equation}\label{eq:uM}
  U^M = u\!\left(\frac{c}{s_M\, e(k)}\right) - \phi_w(\ell_w, k) - \phi_h(\ell_h) - \kappa_S\, S,
  \qquad u(x) = \frac{x^{1-\gamma} - 1}{1-\gamma},
\end{equation}
and the separated wife's is $U^S = u(c/e(k)) - \phi_w(\ell_w, k) - \kappa_S S$. Spending $c$ is divided by the
couple's equivalence scale $s_M$ and the children's scale $e(k)$ (McClements); risk aversion is $\gamma = 1.5$; the
monthly discount factor is $\beta = 0.96^{1/12}$, and the precautionary motive is disciplined by the estimated borrowing
limit rather than by $\beta$. The wife's disutility of work is
\[
  \phi_w(\ell, k) = \mathbf{1}\{\ell = F\}\,\phi_F + \mathbf{1}\{\ell = P\}\,\phi_P + \mathbf{1}\{\ell > 0\}\,\theta(k),
  \qquad \theta(0) = 0,\; \theta(1) = \theta_{05},\; \theta(2) = \theta_{617},
\]
where $k \in \{0, 1, 2\}$ means no child under 18, a youngest child aged 0--5 and a youngest child aged 6--17; the
dependence on the child's age is what the added-work evidence requires.\footnote{After the husband's displacement,
wives with a youngest child under six work 112 more hours a year (s.e.\ 31) and are 4.1 points more likely to be
employed (s.e.\ 0.018); with a youngest child of school age their hours do not rise ($-26$, s.e.\ 36) and their
employment falls by 3.8 points (s.e.\ 0.018) (Section~\ref{sec:household}; \texttt{IGE/19}, \texttt{IGE/38}).} The
husband's disutility is $\phi_h(F) = \phi_h$ and $\phi_h(P) = 0$; his hours choice exists so that the estimator of
\citet{blundell2016consumption}, identified from the intensive margin of both spouses, does not return zero by
construction on simulated data (Section~\ref{sec:validation}). Each employed spouse costs the household a fixed cost
of work $F$ a month, in the budget constraint: it produces the complementarity of consumption and the wife's
participation that \citet{blundell2016consumption} attribute to fixed costs (their Table 8), and it means that a joint
loss lowers spending through income and through the end of two work-related expenditure flows.\footnote{Spousal
leisure complementarity, their other nonseparability, is not modelled and is reported as a known difference.} Taking
up food stamps costs $\kappa_S$ in utility, so that take-up among eligible households is incomplete and rises after a
job loss.\footnote{Food stamp receipt rises by 4.2 points (Black) and 3.3 points (White) in years 1--3 after the
father's displacement, while AFDC/TANF receipt barely moves (Section~\ref{sec:household}; \texttt{IGE/19}).} Each
discrete alternative carries an i.i.d.\ type-I extreme value taste shock with scale $\sigma_\varepsilon$.

\subsection{Wages and human capital}

Spouse $j \in \{w, h\}$ has a human capital level $h_j \in \{1, \dots, 7\}$ and an hourly wage
\begin{equation}\label{eq:wage}
  \ln w_{jt} = \mu_j(r) + \psi_j\,(h_{jt} - 4) + \nu_{jt},
\end{equation}
where $\mu_j(r)$ is a race-specific level, $\psi_j$ the log-wage step between adjacent levels and $\nu_{jt}$ a
transitory shock. Monthly earnings are $y_j = w_j H(\ell_j)$ with $H(F) = 173$ and $H(P) = 87$ hours. There is no
separate age profile: wage growth over the life cycle comes from human capital accumulation at work, and the bounded
grid flattens it at older ages. The level changes at the end of the month in three steps, truncated to the grid. (i)
\emph{Accumulation and depreciation}: a spouse who worked and was not displaced rises one level with probability
\[
  \pi^{up}_j(\ell_j, h_j) = \pi^{up}_j\,(h_j/4)^{-\eta_{up}}\,\big(\mathbf{1}\{\ell_j = F\} + \kappa_{up}\,\mathbf{1}\{\ell_j = P\}\big),
\]
so that accumulation slows as the level rises, the one-parameter form of \citet{bacher2025joint} normalised at the
middle level;\footnote{The PSID two-year log wage growth of continuously employed wives falls from 0.080 at 25--34 to
0.046 at 35--44 and 0.027 at 45--54 (husbands 0.067, 0.029, 0.013; \texttt{IGE/60.PaperATargets.py}). With a constant
accumulation probability the bounded grid alone flattens this too little. At the start values, with $\pi^{up}_j$
rescaled each time to hold the growth at 25--34, $\eta_{up} = 0.5$ lowers the wage-growth block of the objective from
158 to 150, and $\eta_{up} = 1, 1.5, 2, 2.5$ to 145, 136, 123 and 109; at 0.5 the growth at 45--54 is still 0.045 for
wives and 0.037 for husbands, and at 2.5 the husbands' is 0.032 against 0.013. Matching the mean levels by age in
closed form would need $\eta_{up}$ between 6 (wives) and 16 (husbands), outside its bounds $[0, 3]$: the age profile
of accumulation needs more than level dependence, and the remaining gap is reported as a misfit. $\eta_{up}$ is
estimated from a start of 0.5 (\texttt{IGE/paperA\_model/analytic\_start.py --ext hc\_eta};
\texttt{IGE/paperA\_model/results/diagnostics\_20261018.md}, step 4).} a spouse who did not work falls one level with
probability $\pi^{dn}_j$. (ii) \emph{Displacement loss}: a displaced spouse falls by $d_j$ levels in expectation
($\lfloor d_j \rfloor$ levels, one more with probability $d_j - \lfloor d_j \rfloor$).\footnote{The husband's labor
income is 17\% (Black) and 24\% (White) lower in years 1--3 after his displacement and still about \$6,000--7,000 a
year lower six years later, while his hours recover within four to five years (Section~\ref{sec:household}). Whether a
loss at displacement is needed on top of depreciation is decided by the earnings on re-employment by the length of
non-employment after an involuntary separation (SIPP): depreciation alone makes the loss grow with the spell from
close to zero, while $d_j$ adds a loss after even a short spell; if the loss after 1--3 months is close to zero, $d_j$
is cut.} (iii) \emph{Persistent innovations}: with probability $p^c$ a shared innovation moves both spouses' levels
one step in the same direction (up or down with equal probability), and independently with probability
$p^\varepsilon_j$ an idiosyncratic innovation moves spouse $j$'s level. The innovations are the permanent wage shocks of
\citet{blundell2016consumption} on the human capital grid: without truncation the annual variance of spouse $j$'s
permanent component is $12\,\psi_j^2\,(p^c + p^\varepsilon_j)$ and the covariance across spouses
$12\,\psi_w\psi_h\,p^c$, and $(p^c, p^\varepsilon_w, p^\varepsilon_h)$ are set from their Table 3 (permanent variances
0.0382 for wives and 0.0303 for husbands, covariance 0.0027).\footnote{Joint displacement years raise the covariance of
the spouses' earnings residuals in their panel ($+0.193$, $p=0.019$), not of their wage residuals: the joint loss is an
employment shock and is modelled as one, separately from the shared wage shocks (Appendix~\ref{app:bps}).}

\paragraph{Wage shocks.} On top of the human capital level, each spouse's wage has a permanent and a transitory
shock, correlated across spouses as in \citet{blundell2016consumption}, and the transitory shock is drawn once a year and
held for the twelve months of the year. The choice of frequency follows from what each source measures.
\citet{blundell2016consumption} estimate the transitory shock as a serially uncorrelated component of the annual hourly
wage, annual earnings divided by annual hours, from two-year differences of biennial data; they describe it as shocks
such as short illnesses, assume that it is fully smoothed, and estimate variances of 0.0275 for husbands and 0.0125 for
wives with a correlation of 0.31 between spouses, part of which is measurement error from dividing earnings by hours,
which they net out with the shares of a PSID validation study \citep{bound1994evidence}. \citet{bacher2025joint} have no
transitory wage shock: wages are a fixed share of a human capital level that moves up or down one step each month with
some probability, and monthly earnings risk comes from job loss and human capital. High-frequency administrative data
show that monthly earnings within ongoing jobs are volatile but that wages are largely stable from month to month, the
volatility coming from hours set by firms \citep{ganong2026earnings}. Reproducing the annual variance of
\citet{blundell2016consumption} with a shock drawn afresh every month would require a monthly standard deviation of about
0.57, a typical monthly swing in the hourly wage of more than half, which the monthly data rule out. The transitory shock
is therefore an annual shock; part of its variance in \citet{blundell2016consumption} is treated as measurement error,
added only when simulated couples are turned into PSID-style data for the comparison with their estimates; and the
month-to-month risk in the model comes from displacement and job finding. Each spouse's shock takes two values,
$\pm\sigma_{\nu,j}$, and the pair has probabilities $(1+\rho_\nu)/4$ for equal signs and $(1-\rho_\nu)/4$ for opposite
signs, so that the variances are $\sigma_{\nu,j}^2$ and the correlation $\rho_\nu$. The variances start at the
estimates of \citet{blundell2016consumption}, an upper bound, and are reduced once, before estimation, by the part that
their estimator attributes to transitory variance when it is run on the simulated panel with the measurement-error
layer; they are then held fixed.

\subsection{Displacement risk in three parts}

A job ends involuntarily only through a displacement, a permanent job loss.\footnote{Temporary layoffs with recall are
outside the model; they account for about a third of the CPS same-month excess, and the CPS targets use non-temporary
job losses (excess 2.6 points on a base of 0.23\%; Appendix~\ref{app:robust}).} At the end of the month the
displacement shocks are drawn in two stages. With probability $\pi_c(s)$ a common shock hits the couple and every
employed spouse is displaced. Otherwise each employed spouse is displaced independently with probability
$\delta_j(r)\, m_\tau$, where $m_L = 1$ and $m_H > 1$ is the risk multiplier of high-risk couples, whose share is
$p_H$. When both spouses work,
\begin{equation}\label{eq:both}
  \Pr(\text{both displaced}) = \pi_c(s) + \big(1 - \pi_c(s)\big)\, m_\tau^2\, \delta_w(r)\,\delta_h(r),
\end{equation}
and only spouse $j$ is displaced with probability $(1 - \pi_c(s))\, m_\tau \delta_j(r)\,(1 - m_\tau \delta_{-j}(r))$.
The three parts leave different traces at the monthly frequency: the common shock produces excess joint loss only in
the month it hits, the couple type a flat elevation of the wife's hazard in every month around the husband's loss and
in the PSID placebo window, and the own hazards the base rates. The monthly elevation pins the couple type along a
ridge rather than at a point: a rare and very risky type and a common and mildly risky one give the same CPS leads and
lags, but the rare type gives the smaller PSID same-year excess, because its displacements are concentrated in few
couples. The estimation starts on this ridge at $p_H = 0.02$ ($m_H = 15.95$), where the annual excess is close to the
data, and estimates both.\footnote{Along the ridge (own hazards rescaled to hold the base rates), the PSID excess
$\Pr(\text{she displaced} \mid \text{he is}) - \Pr(\cdot \mid \text{he is not})$ falls from 0.118 at $p_H = 0.10$ to
0.071 at $p_H = 0.02$ (data 0.063, s.e.\ 0.005), while the CPS leads and lags stay at 0.0034--0.0061 (data
0.0056--0.0067); off the ridge, a smaller $m_H$ at $p_H = 0.10$ fits the annual excess only by losing the leads and lags
(at $m_H = 4$ they fall to 0.0015). The joint-loss and base-hazard blocks of the objective are 401 at $p_H = 0.10$ and
185 at 0.02, so the annual moments separate $p_H$ along the ridge. The alternative tested and rejected is a short memory
of the common shock: each spouse's hazard multiplied by $1 + c_{mem}$ while the other spouse is non-employed and
entitled after a displacement, one more parameter. At its start ($c_{mem} = 5.82$, $m_H = 4.33$) the joint-loss blocks
are 177 against 62 on the ridge, and the excess one month before and after his loss is 0.013--0.014 against 0.006 in
the data. It stays in the code as an option (\texttt{estimate.py --cs-memory on}). Code:
\texttt{IGE/paperA\_model/results/diagnostics\_20261018/profile\_sigma\_ridge.py}; numbers in
\texttt{IGE/paperA\_model/results/diagnostics\_20261018.md}, B and steps 2a and 2.} The model has no aggregate state;
the larger excess in recessions (6.7 against 3.6 points) is used in a comparative static.

\subsection{Search friction and unemployment insurance}

A spouse who is non-employed at the start of a month finds a job by its end with probability $\lambda_j(r)$, which
measures the search friction; it depends on race and spouse only, not on age, the human capital level or the length
of the spell.\footnote{In \citet{bacher2025joint}, where job finding is set by firms, the rate for non-employed spouses
falls from about 0.16 a month at 26 to about 0.01 at 64 (their Figure 10), mostly through the age distribution of
assets and human capital. Job finding by age is a validation (Section~\ref{sec:validation}); one age slope in
$\lambda_j$ is added only if it fails.} A job allows either hours level. A wife without entitlement to unemployment
insurance who finds one chooses at the start of the next month whether and how much to work, and declining it returns
her to non-employment. A wife who is entitled decides at the end of the month whether to accept, as in
\citet{bacher2025joint}: declining keeps her entitlement, which then lapses at its monthly hazard. The husband always
takes a job he finds. A wife can also leave a job she holds; her monthly exit from employment outside a job loss is a
target, so that margin is disciplined by data.

A displaced spouse becomes entitled to unemployment insurance with probability $q_{UI}$, which stands for take-up and
monetary eligibility together; a spouse who leaves a job is not entitled. An entitled spouse who is still non-employed
loses entitlement at the end of each month with probability $1/6$, so that the expected duration is the 26 weeks of
regular benefits; this geometric expiry is that of \citet{bacher2025joint} and replaces a duration clock, which the
moments that see benefit duration cannot distinguish from it. While entitled, spouse $j$ receives
\[
  b_j = \min\big\{\rho_{UI}\; \bar w_j(h_j)\, H(\ell_j^{lost}),\; \bar b_{e}\big\},
  \qquad \bar w_j(h) = \exp\{\mu_j(r) + \psi_j(h - 4)\},
\]
with replacement rate $\rho_{UI} = 0.5$, a cap $\bar b_e$ equal to the mean of the states' maximum weekly benefits in
policy era $e$, and $\ell_j^{lost}$ the hours of the lost job, which the wife's entitlement state records and which is
full time for the husband. The wife's labor state therefore takes four values (employed; non-employed and entitled
after a part-time job; after a full-time job; not entitled) and the husband's three.

\subsection{The government and the budget constraint}

The government is represented by its schedules, because the paper's policy question is what each schedule pays in
the state in which no one earns.\footnote{The tax-and-transfer function of \citet{blundell2016consumption} is
log-linear in positive family earnings and contains the federal income tax, the EITC and food stamps, but no
unemployment insurance; the state in which no one earns is outside its domain.} All schedules are set by policy era
$e \in \{1984\text{--}96,\, 1997\text{--}2008,\, 2009\text{--}21\}$, each solved with its own stationary schedule,
couples re-optimising at the boundaries. On positive annualised income $Y^a = 12\,(y_w + y_h + b_w + b_h)$ the
household keeps $(1 - \chi_e)\,(Y^a)^{1 - \mu_e}$, the functional form of \citet{blundell2016consumption}, with
$(\chi_e, \mu_e)$ fitted by era to the federal income and payroll taxes of married couples filing jointly, without the
EITC and food stamps, which enter separately; the monthly tax is $T_e = [Y^a - (1-\chi_e)(Y^a)^{1-\mu_e}]/12$. The
Earned Income Tax Credit is paid monthly on annualised family earnings $Y = 12\,(y_w + y_h)$,
\begin{equation}\label{eq:eitc}
  \mathcal E = \tfrac{1}{12}\,\mathrm{EITC}_e(Y; n), \qquad
  \mathrm{EITC}_e(Y; n) = \max\Big\{0,\; \min\{\epsilon^{in}_{e,n} Y,\, \bar E_{e,n}\} - \epsilon^{out}_{e,n}\max\{0,\, Y - Y^{out}_{e,n}\}\Big\},
\end{equation}
with the phase-in rate, maximum credit, phase-out rate and threshold for married filers with $n$ children; it is zero
when no one earns. A household of size $n_{hh}$ is eligible for food stamps when gross income
$Y^g = y_w + y_h + b_w + b_h$ is at most 130\% of the poverty line and assets are below the era's limit; its benefit is
the maximum allotment minus 30\% of net income after the 20\% earned-income deduction, the standard deduction and the
dependent-care deduction,
\[
  B_S = \max\big\{0,\; A_e(n_{hh}) - 0.3\,\max\{0,\; Y^g - 0.2\,(y_w + y_h) - D_e(n_{hh}) - c_{cc}\}\big\},
\]
and take-up $S \in \{0, 1\}$ is chosen. The wife's work costs childcare $c_{cc} = p_{cc}(k)\,H(\ell_w)$ at an hourly
price that depends on the youngest child's age; a subsidy lowers $p_{cc}$. The number of children $n$ is not a state:
it is the PSID mean for couples with the same youngest-child bucket, wife's age group and race. A transfer $Tr \ge 0$
guarantees equivalised spending of at least $c_{min}$ (the AFDC/TANF maximum for a family of three in the era),
standing for the transfers that are not modelled. Cash on hand of a married couple is
\begin{equation}\label{eq:budget}
  \begin{aligned}
  X^M = {} & \big(1 + r_m + \omega\,\mathbf{1}\{a < 0\}\big)\,a + y_w + y_h + b_w + b_h - T_e + \mathcal E \\
           & + S\, B_S - c_{cc} - F\,\big(\mathbf{1}\{\ell_w > 0\} + \mathbf{1}\{\ell_h > 0\}\big) + Tr,
  \end{aligned}
\end{equation}
and the household chooses $c$ and next month's assets $a' = X^M - c \ge -\underline{a}$. Saving earns
$r_m = 1.03^{1/12} - 1$; debt costs the real interest rate on credit cards, 11.4\% a year, so that the wedge $\omega$ is
8.4 points a year, a coded input and not a parameter.\footnote{Federal Reserve G.19, interest rate on credit card plans
of commercial banks, all accounts, 13.7\% nominal on average over the 108 quarters of 1995--2021, less 12-month CPI-U
inflation (2.3\%), less the real saving rate of 3\%; the rate on accounts assessed interest gives 12.2\% real. The
liquid wealth of the targets nets unsecured debt against liquid assets, and couples with negative liquid wealth are
mostly debtors (\texttt{ReducedForm/62.PoorMarriedCouples.py}), so the credit-card rate is the marginal rate on that
debt (\texttt{IGE/paperA\_model/borrowing\_wedge.py}, \texttt{inputs/borrowing\_wedge.json}). At the start values and
on common draws the wedge lowers the objective from 68{,}086 to 64{,}657 and the share of wives aged 25--34 in couples
with liquid wealth at or below zero from 0.80 to 0.48 (data 0.57); it lowers the objective at every borrowing limit
tried (63.5--65.6 thousand against 66.5--69.8 thousand without it), and with the wedge the wealth block falls from 3{,}079
to 2{,}672, 2{,}419 and 2{,}317 as $\underline a$ rises from 15 to 25.6, 40 and 55 thousand dollars. The shape at 45--59
stays wrong with or without the wedge: too few couples at or below zero (0.085 against 0.381) and too many above one
year of earnings (0.388 against 0.177). It is reported as an open misfit
(\texttt{IGE/paperA\_model/results/diagnostics\_20261018.md}, step 3). Part of it is the definition of the target: the
model's single asset also funds the retirement terminal value, while liquid wealth leaves retirement saving out.
Measured as net worth without home equity and vehicles, on the same couples and waves, the share above one year of
earnings at 45--59 is 0.399 (model 0.388) and the shares at or below zero at 25--34 and 35--44 are 0.476 and 0.354
(model 0.476 and 0.340); what remains is the share at or below zero at 45--59, 0.264 against 0.085, couples who save
nothing (definition minus liquid at 45--59: $-0.117^{***}$ and $+0.222^{***}$, both $p<10^{-15}$;
\texttt{IGE/77.WealthDefinitions.py}).} The separated wife's cash on hand is the same
expression without the husband's terms, with the single-filer EITC and food-stamp schedules.

\subsection{Marriage and separation}

The marriage ends at the end of a month with probability
\begin{equation}\label{eq:sep}
  \sigma(r, e_h') = \sigma_0(r) + \sigma_J(r)\, \mathbf{1}\{e_h' = 0\},
\end{equation}
which is higher while the husband is out of work after a displacement ($e_h' = 0$ is his status at the start of the
next month; he is non-employed only after a displacement). The hazard responds only to his displacement, whereas in
the PSID her displacement is followed by break-up as often as his (Section~\ref{sec:household}); a hazard that also
responds to her displacement is a candidate extension after the first estimate. At a separation assets are divided equally and the wife
continues alone with the children, with her own labor state, human capital, entitlement, transitory shock and child
state; there is no remarriage within the horizon and no child support.\footnote{In the PSID, separated mothers with
children who remain single receive on average \$132 (Black) and \$327 (White) a month in child support (2000 dollars;
38\% and 64\% receive any; difference $-196^{***}$, $p<0.001$), and remarry at annual rates of 2.6\% and 6.1\%
($-3.5$ points$^{***}$, $p<0.001$). Both would make the separated state less bleak, more so for White wives; they are
inputs without parameters, to be added after the first estimate. Separated mothers work more than married mothers
(within mother $+3.5^{**}$, $p=0.016$, and $+11.8^{***}$ points), an untargeted check on the separated state
(\texttt{IGE/76.SeparatedMothers.py}).} The separation is exogenous: the model
reproduces the separation response of Section~\ref{sec:household} and its difference by race, and the household
anticipates the hazard, which adds a precautionary motive.\footnote{The share of displaced fathers still living with
the wife they had before the job loss falls by 4.5 points (Black) and 1.9 points (White) over years 1--3
(\texttt{IGE/19}). Tying the hazard to his non-employment adds no state variable; a twelve-month flag after his
displacement and a hazard that also responds to her displacement were considered and not adopted for the first
estimate. The second is coded as an option that is off in the first estimate (a term $\sigma_{Jw}$ while she is
entitled after her own displacement, \texttt{--her-sep}): at the start values it moves the simulated response to her
displacement from $-0.001$ to $-0.019$ against $-0.050$ in the data (\texttt{IGE/64}, \texttt{IGE/65}) and lowers the
objective on the extended set of 89 moments from 13{,}232.5 to 13{,}162.8; it is judged at the first estimate
(\texttt{IGE/paperA\_model/results/extensions\_20261019.md}).} The hazard is a
reduced form for the marital surplus, and the marriage is kept exogenous because the data pin the mean response and
little that would discipline a decision: the response is the same after plant closures as after layoffs, which is
what a fall in household income and surplus predicts rather than news about the husband \citep{charles2004job}, and
it does not differ significantly by race, by the child's age or by the wife's employment.\footnote{Still living with
the baseline wife, years 1--3 after his displacement: closures minus layoffs 0.000 ($p = 0.99$, Black) and 0.006
($p = 0.75$, White); by the youngest child's age $-0.030$ and $-0.030$; wife employed at baseline $-0.036^{***}$
against not employed $-0.019$ (difference 0.017, s.e.\ about 0.017); the race difference in the targeted mean is 0.5
standard errors. The outcome is one by construction before the displacement, so it has no placebo test
(\texttt{IGE/19.FatherDisplacement.py}, \texttt{IGE/38.PSIDJointLossChecks.py}). The data's gap keeps widening after
he is re-employed: $-0.027$, $-0.037$ and $-0.050$ in years 1, 3 and 6 (Black $-0.037$, $-0.053$, $-0.074$; White
$-0.017$, $-0.022$, $-0.030$).} Its time path is an untargeted test of the form: the hazard is raised only while he is
out of work, whereas the PSID gap keeps widening long after he is back at work. At the start values the simulated
gap is $-0.029$, $-0.029$ and $-0.039$ in years 1, 3 and 6 against $-0.027$, $-0.037$ and $-0.050$ in the PSID; if the
estimate leaves the same flattening, the alternative is a hazard that depends on his human capital level after he is
re-employed, which adds one parameter and no state. The husband's welfare after a
separation is not counted.

\subsection{Children}

The children block is outside the estimation: the parents' model is estimated without it, so its results for
consumption and the ranking of policies stand on their own. For each child in a simulated household it maps family
income by childhood age and the presence of both parents into an adult outcome,
\begin{equation}\label{eq:child}
  \ln Y^{adult}_i = \alpha + \beta_{05} \ln \bar y_{i,0\text{--}5} + \beta_{617} \ln \bar y_{i,6\text{--}17} + \beta_P\, P_i + \eta_i,
\end{equation}
where $\bar y_{i,g}$ is mean family income (earnings, unemployment insurance, EITC and food stamps) over the months in
which the child is in age group $g$ and $P_i$ the share of those months with both parents present. The slopes are
not estimated. Two scenarios bracket them: a high scenario in which a father's displacement during childhood, as
simulated, lowers adult earnings by about 9\%, the estimate in Canadian administrative data
\citep{oreopoulos2008intergenerational} and the PSID association of \citet{page2009parental}, and a low scenario at the
smaller effects in US tax data \citep{hilger2016parental}; the PSID associations of Section~\ref{sec:children} are an
upper bound only. The presence term is kept because the separation of the couple is one of the channels by which the
joint loss reaches the children; $\beta_P$ is the PSID association in the high scenario and zero in the low one.

\subsection{The household's problem}

The married couple's state at the start of month $t$ is
\[
  \Omega_t = \big(t,\; a,\; e_w, u_w,\; h_w,\; e_h, u_h,\; h_h,\; k,\; \nu;\; s, \tau, r\big),
\]
with assets $a$ (30 grid points, dense near the borrowing limit), employment $e_j \in \{0, 1\}$ and entitlement $u_j$
(the wife's recording the hours of the lost job), the two human capital levels, the youngest-child bucket $k$, this
year's transitory shocks $\nu = (\nu_w, \nu_h)$ (four points), and the permanent characteristics: 211,680 points a
month for each of the twelve problems $(r, s, \tau)$, which are solved and simulated one at a time. The separated
wife's state is $\Omega^S_t = (t, a, e_w, u_w, h_w, k, \nu_w; \tau, r)$. Let $\mathcal D(\Omega, \nu)$ be the feasible
alternatives. The value at the start of the month, after $\nu$ is observed and before the taste shocks, is
\begin{equation}\label{eq:VM}
  V^M_t(\Omega, \nu) = \sigma_\varepsilon \ln \sum_{d \in \mathcal D(\Omega,\nu)} \exp\!\big(v^M_t(d; \Omega, \nu)/\sigma_\varepsilon\big) + \sigma_\varepsilon\gamma_E ,
\end{equation}
with $\gamma_E$ Euler's constant; the choice-specific value is
\begin{equation}\label{eq:vM}
  v^M_t(d; \Omega, \nu) = \max_{a' \ge -\underline a} \Big\{ u\!\Big(\frac{X^M(d; \Omega, \nu) - a'}{s_M\, e(k)}\Big)
     - \phi_w(\ell_w, k) - \phi_h(\ell_h) - \kappa_S S + \beta\, W^M_{t+1}(a'; d, \Omega, \nu) \Big\},
\end{equation}
and the continuation value integrates over next month's state and the separation,
\begin{align}\label{eq:WM}
  W^M_{t+1}(a'; d, \Omega, \nu) = \sum_{\Omega'} P(\Omega' \mid \Omega, d, a')
   \Big[ &\big(1 - \sigma(r, e_h')\big)\, \E_{\nu' \mid \nu} V^M_{t+1}(\Omega', \nu') \nonumber \\
         &+ \sigma(r, e_h')\, \E_{\nu_w' \mid \nu_w} V^S_{t+1}\big(\Omega'^S, \nu_w'\big) \Big],
\end{align}
where $\Omega'^S$ keeps the wife's components of $\Omega'$ with assets $a'/2$. The expectation over the shocks is
degenerate within the year and is the four-point sum over a new draw when $t+1$ is a January. The transition $P$
factors into the events of the month---displacement (the common shock, then the own hazards, depending on $s$, $\tau$,
$r$ and on who worked), job finding, the wife's quit, the three human capital steps, the entitlement states and the
child bucket---and, conditional on the month's displacement and job-finding events, the remaining transitions are
independent, so the expectation is taken one dimension at a time. The separated wife's problem has the same form with
$d = (\ell_w, S)$, cash on hand $X^S$, the scale $e(k)$, her own hazard $\delta_w(r)\,m_\tau$ (the common shock no
longer applies), job finding $\lambda_w(r)$ and a continuation value over her own state only.

The bucket $k$ follows a monthly Markov chain estimated on PSID couples: a birth moves $k = 0$ to $k = 1$ with
probability $b_0(\text{age}_w, r)$ and $k = 2$ back to $k = 1$ with probability $b_2(\text{age}_w, r)$; the youngest
child ages from $k = 1$ to $k = 2$ with probability $1/72$ and leaves childhood from $k = 2$ with probability $1/144$.
At the wife's age 60 the couple retires: the terminal value is the discounted utility of 240 months of spending equal
to the annuity value of $a$ plus Social Security benefits of 40\% of each spouse's earnings at its final level. Each
simulated couple is a copy of a PSID couple, entering in January of the year of its first observation from 1984 on with
its calendar year, the wife's age, race, workplace match, liquid wealth, each spouse's employment, the youngest
child's bucket and each spouse's human capital level, set from the log hourly wage at entry; the couple type is drawn
with probability $p_H$ independently of the observables. Each PSID couple is copied several times, and all draws are
common random numbers across evaluations and counterfactuals. The model is solved by backward induction over 420
months for each problem and policy era; within a month the discrete alternatives are handled with the extreme-value
shocks and consumption with the endogenous grid method for discrete-continuous models \citep{carroll2006method,
iskhakov2017endogenous}, as in \citet{bacher2025joint}. Between asset nodes the continuation value is a
monotone cubic Hermite interpolant of the node values and their exact derivatives. The food-stamp asset test puts a
notch in next month's value at the asset limit $L$ (1{,}884 dollars in the 1997--2008 schedule): a couple that saves
beyond it loses the food-stamp alternative. The solver therefore carries grid nodes at $L$ and just above it, treats
the corner $a' = L$ as a candidate in every state, and runs a two-branch endogenous grid where the grid folds at the
notch.\footnote{Without this treatment the interpolated continuation value smeared the drop and the first-order
conditions missed the bunching at $a' = L$: the gap between the endogenous-grid solution and a brute-force maximisation
over $a'$ with the same continuation values reached 0.16 utility units (99th percentile $3.4 \times 10^{-3}$) at low
cash on hand. With it the 99th percentile and the maximum of the gap are $1.7 \times 10^{-4}$ and $2.6 \times
10^{-3}$ at the start values, $7.0 \times 10^{-4}$ and $2.4 \times 10^{-3}$ with $\kappa_S = 0$, and $3.7 \times
10^{-5}$ and $1.2 \times 10^{-4}$ with $\kappa_S = 2.64$, against criteria of $10^{-3}$ and $5 \times 10^{-3}$; in 79\% of
states the endogenous grid is monotone, or monotone on both sides of $L$. All twelve staged checks pass (\texttt{IGE/paperA\_model/tests/run\_all.py};
\texttt{IGE/paperA\_model/results/diagnostics\_20261018.md}, follow-up 1).} The simulator writes monthly histories, and each data set is
reproduced by writing the histories in that data set's shape---PSID household-years in the waves in which the couple is
observed, CPS couple-months in the four-eight-four rotation, SIPP monthly paths, the household-year panel of
\citet{blundell2016consumption}---and running the data scripts' definitions on it.

\subsection{Estimation}\label{sec:estimation}

All endogenous parameters are estimated jointly in one simulated-method-of-moments run; nothing is tuned by hand, and
every calibrated input is a coded mapping from data or an external source: the preferences shared with the lifecycle
model of women's careers that shares this code ($\beta$, $\gamma$, the equivalence scales, the hours levels,
$\kappa_{up} = 0.5$, and the taste-shock scale $\sigma_\varepsilon = 0.05$, which is a fixed scale and not
estimated\footnote{A larger taste shock does not move the wife's exit rate by the order of magnitude the data need. At
the start values (one-month CPS flows, before the extensions) $\sigma_\varepsilon = 0.05$, 0.10, 0.20 and 0.30 give a
monthly exit rate of 0.26\%, 0.35\%, 0.29\% and 0.20\% against 2.6\%, while the employment of wives with no child under
18 rises from 0.95 to 0.97, 0.99 and 0.99 (data 0.81), the entry rate of White wives from 0.043 to 0.113 (data 0.075),
and the added-worker response in employment falls from 0.006 to $-0.003$ (data 0.042): the option value of holding a
job grows with $\sigma_\varepsilon$, so wives keep their jobs. The objective is 50{,}920, 44{,}860, 50{,}652 and
57{,}938 (\texttt{IGE/paperA\_model/results/diagnostics\_20261018/profile\_sigma\_ridge.py}; the code can free it,
\texttt{estimate.py --free-extra sigma\_eps}).}); the borrowing wedge (G.19); the wage levels and steps (PSID log hourly wages by race; seven levels
spanning the 1st--99th percentiles of the residual at 40--49); the wage-shock probabilities and transitory variances
(Table 3 of \citealp{blundell2016consumption}); the workplace-match shares and the child process (PSID couples); the
schedules by era (the unemployment-insurance cap from the states' maximum weekly benefits,
\texttt{IGE/12b.UIGenerosity.py}; take-up $q_{UI}$ from the Displaced Worker Supplement; EITC, food-stamp and tax
parameters from the IRS, the USDA and TAXSIM); the childcare price (PSID childcare spending per hour of the mother's
work by child age); and the consumption floor. Table~\ref{tab:model_params} lists the thirty-two estimated parameters
with the moments that identify them. Estimation runs at a reduced size: seven human capital levels, the shared-employer
and same-industry groups merged into one group that shares $\pi_c(1)$, the 1997--2008 schedule in every year, and 20
asset points, so that an evaluation solves eight problems instead of thirty-six; the counterfactuals run at the full
size. With the two groups merged $\pi_c(2)$ does not enter, and thirty-one parameters are free. The cut list, in
order, is a common $\delta_w$ across races, $\phi_h$ calibrated to the husbands' part-time share, a common $\pi^{dn}$
for both spouses, $d_j$ if the re-employment earnings by spell length are fit by depreciation alone, and the wife's
lost-hours record in the entitlement state.

\paragraph{Targets.} The targets are the moments of \texttt{IGE/60.PaperATargets.py}; the first evaluations of the
coded model changed the definition of the CPS flows and the set of moments in the objective. At the start values the model's wives left employment at 0.26\% a month against 2.6\% in the
CPS, and a larger taste shock did not close the gap. Much of the CPS flow consists of one-month reversals, a wife
recorded out of work for a single month between months of work, which the model's monthly statuses do not produce. The
entry and exit rates of wives are therefore computed with one filter in the data and in the simulated CPS panel: a
transition counts only if the old status held in months $t-1$ and $t$ and the new status holds in $t+1$ and $t+2$, all
four months in one rotation spell, with the husband employed in $t$ and $t+1$. The filter cuts the data's exit rate
from 0.0261 to 0.0101 and the entry rates from 0.0751 to 0.0320 (White) and from 0.0942 to 0.0404 (Black); the model's
flows are almost all persistent already.\footnote{At the extended start the model's one-month and persistent rates are
0.0023 and 0.0014 (exit), 0.047 and 0.036 (White entry), 0.099 and 0.089 (Black entry); the objective falls from 90{,}296
to 12{,}761 with the persistent flows (\texttt{IGE/60.PaperATargets.py}; \texttt{IGE/paperA\_model/panels.py};
\texttt{IGE/paperA\_model/results/diagnostics\_20261018.md}, follow-up 2). The correction for response errors of
\citet{poterba1986reporting} was coded first and dropped (\texttt{estimate.py --cps-persist off --cps-classif on}):
applying their error rates for women aged 25--59 to the simulated monthly statuses brings the exit rate to 0.0240
against 0.0261, but the entry rates to 0.208 and 0.364 against 0.075 and 0.094. At the data's employment rate of
0.81 the spurious entries of misrecorded employed wives alone give a measured entry rate of at least 0.083, above the
White target, so the true entry rate that their errors imply is negative. Their errors are for 1981 and serially
uncorrelated, which they note overstates the adjustment, and the CPS targets start in 1994, after the redesign with
dependent interviewing.} Even so, the persistent exit rate is the largest misfit at the start (0.0014 against 0.0101;
4{,}623 of an objective of 11{,}734): wives in the model leave jobs about seven times less often than in the
data.\footnote{The objective at the start is 11{,}733.7 from a clean compilation on the PC and 11{,}733.6 on the
cluster (job 11863423); compiled binaries differ in the fourth significant digit (cached and fresh builds by up to
$9\times10^{-5}$ in a moment), so the reproducibility check compares the objective at a relative $10^{-3}$
(\texttt{IGE/paperA\_model/hpc\_job\_a.py}).} The missing exits are those of wives who have held their job for a
year or more: by job tenure in the SIPP the model's monthly exit is 0.035, 0.010 and 0.0027 at 0--5, 6--11 and 12 or
more months against 0.021, 0.016 and 0.0068 in the data, so a hazard specific to the first year of a job, which raises
the first cell further, is not the remedy.\footnote{Data: \texttt{IGE/66.PaperATargetsExt.py} from the SIPP
2004--2014 panels (household bootstrap standard errors 0.0007, 0.0007 and 0.0001). An extra exit hazard in the first
year of a job is coded as an option that is off in the first estimate (\texttt{--new-job}); at the start values it
raises the first cell to 0.042 and the objective on the extended set of 89 moments from 13{,}232.5 to 13{,}245.7
(\texttt{IGE/paperA\_model/results/extensions\_20261019.md}). Why established wives leave is studied in
\texttt{proposals/wife\_exit\_evidence.md}.}

Of the 97 targets, 85 enter the objective. The other twelve stay in the fit table, one reason each: the four pooled CPS
rates (in the model each is exactly the race-share average of its Black and White rates), the three one-month CPS flows
(replaced by the persistent ones), the three SIPP employment rates of husbands without a separation (in the model
husbands never quit), and the husband's employment rate by race. The model's husbands have positive hours in 99.9\%
(Black) and 100\% (White) of couple-years against 88.2\% and 94.3\% in the PSID: they leave work only by displacement and
find a job in about four months, so a year without work almost never occurs, and no parameter reaches the
target.\footnote{As a target it would add 1{,}036 to the objective ($z$ of 22--23) and pull $\lambda_h(r)$ down against
the SIPP re-employment path; reported only, the objective at the start falls by 1{,}036 to 11{,}734
(\texttt{moments.OBJ\_EXCLUDE}; \texttt{IGE/paperA\_model/results/diagnostics\_20261018.md}, step 5 and follow-up 3).}
The weights are $1/(\mathrm{se}_m^2 + \mathrm{var}^{sim}_m)$, the bootstrap variance of each data moment plus the
simulation variance of the model moment across five sets of common random numbers, so that a moment the simulation
reproduces noisily does not dominate.\footnote{Under inverse bootstrap variances alone the CPS hazards held 96\% of the
weight at the start (91\% with the trim and the simulation variance); the simulation variances are those of the start
before the extensions and are recomputed at the Levenberg--Marquardt candidate before the global search
(\texttt{IGE/paperA\_model/moments.py}, \texttt{weights}; \texttt{inputs/moment\_noise.csv}).}

\paragraph{Procedure.} Every step uses the same common random numbers, three copies of each PSID couple. The start
values are closed forms from the targets and inputs (\texttt{analytic\_start.py}); a Jacobian of the moments at the
start checks that every parameter moves a named block. A Levenberg--Marquardt search in logistic units, with a
central-difference Jacobian at each iteration, takes the start to a candidate; the simulation variances are recomputed
there; and CMA-ES searches around the candidate with initial step sizes from the Jacobian's column norms, from three
seeds, as the check against local minima. The route is first run on pseudo-data generated at known parameters, and the
independent-shock model is estimated by the same route.\footnote{CMA-ES alone was tried first: a pseudo-data recovery
of 19 parameters from 496 evaluations lowered the objective from 2{,}907 to 482 (zero at the truth) and was still
falling, with 3 of 19 parameters within two simulation standard errors of the truth. The jobs on the university
cluster are \texttt{Joint-test}, \texttt{Joint-recovery-gn}, \texttt{Joint-gn-main}, \texttt{Joint-noise} and
\texttt{Joint-smm-main} with two restarts (\texttt{IGE/paperA\_model/estimate.py}, commands \texttt{polish},
\texttt{noise}, \texttt{run}; \texttt{hpc\_job\_a.py}).} Standard errors are of the sandwich form, with the bootstrap
covariance of the targets and the simulation covariance of the moments.

\begin{table}[htbp]\centering\footnotesize
\caption{Estimated parameters and identifying moments}\label{tab:model_params}
\begin{tabular}{p{4.0cm}p{2.6cm}p{8.2cm}}
\toprule
Parameter & Symbol & Identifying moments (window) \\
\midrule
\multicolumn{3}{l}{\textit{Displacement risk (9)}} \\
Wife's own hazard, by race & $\delta_w(B), \delta_w(W)$ & her monthly non-temporary job-loss rate when he keeps his job (CPS 1994--2023); her annual displacement rate (PSID 1984--2021) \\
Husband's own hazard, by race & $\delta_h(B), \delta_h(W)$ & his monthly job-loss rate (CPS); his annual displacement rate, 8.0\% / 4.8\% (PSID) \\
Common shock, by workplace match & $\pi_c(0), \pi_c(1), \pi_c(2)$ & same-month excess by match (CPS); same-year joint displacement by sector (PSID) \\
Couple type & $p_H$, $m_H$ & her excess in months $\pm 1$ to $\pm 3$ around his loss (CPS); permanent component and placebo window (PSID) \\
\multicolumn{3}{l}{\textit{Search friction (4)}} \\
Wife's job finding, by race & $\lambda_w(B), \lambda_w(W)$ & persistent monthly employment entry of non-employed wives (CPS); her employment path 0--6 months after his separation (SIPP 2003--16) \\
Husband's job finding, by race & $\lambda_h(B), \lambda_h(W)$ & his employment 3, 6, 12 months after an involuntary separation (SIPP); his hours in years 0--1 after displacement (PSID) \\
\multicolumn{3}{l}{\textit{Human capital (7)}} \\
Accumulation & $\pi^{up}_w, \pi^{up}_h$ & two-year log wage growth of the continuously employed at 25--34 (PSID) \\
Accumulation by level & $\eta_{up}$ & the fall of that growth at 35--44 and 45--54 (PSID) \\
Depreciation & $\pi^{dn}_w, \pi^{dn}_h$ & earnings on re-employment by length of non-employment, 1--3, 4--12, 13--24 months (SIPP); her wage at re-entry after 12+ months out; his loss in year 6 relative to years 1--3 (PSID) \\
Displacement loss & $d_w, d_h$ & labor income loss in years 1--3 after own displacement (PSID); loss on re-employment after 1--3 months (SIPP) \\
\multicolumn{3}{l}{\textit{Work, safety net, marriage, credit (12)}} \\
Wife's work costs & $\phi_F, \phi_P$ & employment of wives with no child under 18; part-time share (PSID); her persistent monthly exit from employment outside job loss (CPS) \\
Child costs & $\theta_{05}, \theta_{617}$ & her employment by youngest child's age; her added hours and employment after his displacement by child age (PSID) \\
Husband's full-time cost & $\phi_h$ & part-time share of employed husbands (PSID) \\
Fixed cost of work & $F$ & consumption change after her own displacement; joint minus single (PSID 1999--2021) \\
Food-stamp stigma & $\kappa_S$ & receipt of couples; its rise after his displacement, 4.2 / 3.3 points (PSID) \\
Separation, by race & $\sigma_0(B), \sigma_0(W)$ & annual separation rate of couples (PSID) \\
Separation response, by race & $\sigma_J(B), \sigma_J(W)$ & still with wife, $-4.5$ / $-1.9$ points in years 1--3 (PSID) \\
Borrowing limit & $\underline a$ & share of couples with liquid wealth $\le 0$ by wife's age; consumption change after his displacement (PSID 1999--2021) \\
\bottomrule
\end{tabular}
\begin{minipage}{0.95\textwidth}\footnotesize Thirty-two parameters; thirty-one are free at the estimation size, at which
the groups $s = 1$ and $s = 2$ are merged and share $\pi_c(1)$, so that $\pi_c(2)$ does not enter. Every moment is computed on simulated households of
the same window as its data counterpart: PSID couple-years 1984--2021, consumption 1999--2021 only, CPS couple-months
1994--2023, SIPP 2003--2016. Standard errors of the moments are bootstrap standard errors by 1968 family (PSID) or
household (CPS, SIPP); the weights are the inverse of the bootstrap variance plus the simulation variance of each
moment, with equal block weights as a robustness run; 85 of the 97 targets enter the objective. Targets:
\texttt{IGE/60.PaperATargets.py}; consumption: \texttt{IGE/47.PSIDConsumptionEvent.py}.\end{minipage}
\end{table}

\paragraph{Identification.} The monthly frequency separates the three parts of displacement risk. With the husband
employed, the wife's monthly job-loss probability is $m_\tau\,\delta_w(r)$ exactly, and her excess in the months
around his loss is $\delta_w(r)\,\mathrm{var}(m_\tau)/\E(m_\tau)$ to first order: the base rate and the flat excess
give $\delta_w$ and the dispersion of the couple type, the PSID placebo window a second reading of the latter, and the
same-month spike in each workplace group, given these, $\pi_c(s)$, with the PSID same-year moments by sector
over-identifying it. Because the couple type raises both own hazards for years, it is also what makes the risk matter
for saving. The wife's job finding is identified by the entry rate of non-employed wives and by the speed with which
employment builds up after the husband's loss, the work costs by the levels of employment by child age and the
part-time share; the response to his loss separates the two, since a high work cost lowers the eventual response and a
high search friction delays it. Her persistent exits from employment outside job loss were to pin, through the same
work costs, the margin on which she leaves a job; at the start values the work costs alone miss that rate sevenfold,
and the evidence on why wives leave employment decides what the model needs for it. Accumulation is identified by the
wage growth of the continuously employed, its level at 25--34 by $\pi^{up}_j$ and its fall with age by
$\eta_{up}$; the earnings on re-employment by the
length of non-employment separate depreciation (the slope of the loss in the spell) from the displacement loss (its
level after a short spell). The stigma is identified by food-stamp receipt and its rise after displacement, the
separation parameters by the baseline rate and the response, the borrowing limit by the share of couples with no liquid
wealth and the consumption fall after his displacement, and the fixed cost by the consumption change at the wife's own
displacement, which is smaller than her income loss because the cost of work ends. At the start values every
parameter moves a named block of moments; the weakest directions are $m_H$ with $\eta_{up}$, $\pi^{dn}_h$ with $d_h$,
$\lambda_h(B)$ against $\lambda_h(W)$, and $F$, and the Jacobian is read again at the Levenberg--Marquardt
candidate.\footnote{Range-scaled condition number 529 at the extended start, 31 free parameters
(\texttt{IGE/paperA\_model/results/diagnostics\_20261018.md}, C). Ties across races remain an option of the code
(\texttt{--tie}) and are not applied: the racial differences are reported.} A restricted model with
independent shocks, $\pi_c(\cdot) = 0$ and $m_H = 1$, is estimated on the same targets with the same draws and
reported next to the main model in the fit and counterfactual tables; the difference is what the correlation buys.

\subsection{Validation}\label{sec:validation}

The following are computed on the simulated households through the same definitions and windows and are not used in
estimation: the CPS jump in looking the month after his loss (8.8 points) and the employment path of these wives
over twelve months; the SIPP paths from month 12 to month 36; the Displaced Worker Supplement same-year ratio (4.7\%
against 1.1\%); the family income loss after his displacement; the fall in job finding of wives already unemployed
when he loses his job, which the model, in which the common shock does not lower her job finding, is expected to miss;
the excess by child age and by race, equal in the data; assets relative to human wealth against Figure 3 of
\citet{blundell2016consumption}; the level of the PSID age profile of married men's earnings; the monthly transition
matrices of wives and husbands; and job finding by age against the CPS, the SIPP and \citet{bacher2025joint}. Two
validations test channels the model does not parameterise. The added-worker effect of \citet{bacher2025joint}---the
entry of a non-participating spouse in the month the employed spouse loses the job, 5.55 points overall and 7.12,
6.83, 4.57 and 1.40 points from the youngest to the oldest age group---is reproduced on this paper's CPS panel
(\texttt{IGE/50.CPSAddedWorkerByAge.py}) and compared with the simulated effect by the wife's age; since no parameter
depends on age, an age gradient in the model can come only from human capital, assets, the children and the horizon.
And the estimator of \citet{blundell2016consumption} is run, unchanged, on a PSID-shaped biennial panel of about 50,000
simulated couples with calendar years labelled 1998--2008 and measurement error added as a two-component mixture
calibrated to the composition of the rows their trims remove (\texttt{IGE/bps\_reconciliation/}): it must return their
wage variances and covariances within their standard errors, their insurance split within about ten points, their
outside-insurance and extensive-margin coefficients within their standard errors, and the derivatives of these with
respect to their restrictions that Appendix~\ref{app:bps} reports. Because their estimator returned a wife's hours
response to her husband's permanent wage of $-0.76$ to $-0.86$ for a true response of $-0.3$ in a synthetic test, the
test compares the model's own simulated response with what their estimator reports on the model's data, not only the
estimates with theirs. Spousal leisure complementarity and the exact Frisch magnitudes are not expected to match.

\subsection{Counterfactuals}

Each experiment re-solves the model with the same draws, is reported for the main and the independent-shock models,
and is scored on consumption in the month of the husband's loss and over the following twelve months, the share of
couple-months with no earner and with no earner and no credit, a consumption-equivalent welfare change at entry, and
the children's adult outcomes under the high and low slopes, each by workplace match and race. (i) \emph{Independent
displacement at unchanged risk}: how much consumption smoothing and how much of the children's outcomes the correlation
costs. Each spouse draws an own type $\tau$ and an own common-shock draw, so each spouse's monthly hazard
$\pi_c(s) + (1 - \pi_c(s))\,m_\tau\,\delta_j(r)$ keeps its distribution and only the dependence between the spouses
goes; setting $\pi_c(\cdot) = 0$ and $m_H = 1$ instead would also lower each spouse's risk.
(ii) \emph{Unemployment insurance for both spouses and longer in recessions}: both spouses entitled at a joint loss,
and the expiry hazard from $1/6$ to $1/12$ (52 weeks), at the estimated $\pi_c$ and with $\pi_c$ doubled, the
recession excess; the added-worker effect under the longer duration is reported next to Table 12 of
\citet{bacher2025joint}, in which it falls from 5.16 to 3.68 points. (iii) \emph{An earnings subsidy against an
unconditional transfer of equal cost}: an EITC-type expansion against a monthly transfer to couples with children,
scored on consumption in joint-loss months, on months with no earner and no credit, and on children; the subsidy pays
only when someone works, so it insures least in the state in which both spouses have lost their jobs. (iv) \emph{A
childcare subsidy} when the youngest child is under six. (v) \emph{Black couples with White displacement hazards},
then with the White separation response, one at a time. The welfare numbers cover the married household and, after a
separation, the wife's household; the husband is not followed after a separation, a limitation stated with every
welfare number.
